% Generated draft: requires mathematical and source review.
\documentclass[11pt]{article}
\usepackage[margin=25mm]{geometry}
\usepackage{fontspec}
\setmainfont{Latin Modern Roman}
\usepackage{amsmath,amssymb,mathtools,unicode-math}
\setmathfont{Latin Modern Math}
\usepackage{enumitem,needspace,xurl,hyperref}
\setlength{\parindent}{0pt}
\setlength{\parskip}{5pt}
\setlength{\emergencystretch}{3em}
\newcommand{\N}{\mathbb{N}}
\newcommand{\Z}{\mathbb{Z}}
\newcommand{\Q}{\mathbb{Q}}
\newcommand{\R}{\mathbb{R}}
\newcommand{\C}{\mathbb{C}}
\newcommand{\F}{\mathbb{F}}
\newcommand{\A}{\mathbb{A}}
\newcommand{\PP}{\mathbb{P}}
\newcommand{\E}{\mathbb{E}}
\newcommand{\eps}{\varepsilon}
\newcommand{\dd}{\,\mathrm{d}}
\begin{document}
\section*{UnsolvedMath ID: 30002635}
\textbf{Source problem code:} catalogue entry
\subsection{Definitions and mathematical statement}
The supplied record does not specify the dimension, compactness, connectedness, or boundary conditions of \((M,g)\). Consequently, its unrestricted intended scope cannot be determined. The formulation below makes precise the customary closed-manifold interpretation, conditional on confirming the additional hypotheses that \(M\) is a nonempty smooth compact manifold without boundary of dimension \(n\geq 3\), and that \(g\) is a smooth Riemannian metric. These hypotheses are an explicit interpretive completion, not hypotheses stated in the supplied record; connectedness is left unrestricted.

Under this interpretation, let \(G\leq\operatorname{Isom}(M,g)\) be any subgroup. No closedness or finiteness assumption on \(G\) is imposed. A metric \(h\) is \(G\)-invariant if \(\gamma^*h=h\) for every \(\gamma\in G\). Define the \(G\)-invariant conformal class by
\[
[g]_G=\left\{u^{4/(n-2)}g:\ u\in C^\infty(M,\R),\ u>0,\ u\circ\gamma=u\text{ for every }\gamma\in G\right\}.
\]
For a smooth metric \(h\), let \(R_h\) denote its scalar curvature, let \(\dd\mu_h\) denote its Riemannian volume measure, and write \(\operatorname{Vol}_h(M)=\int_M\dd\mu_h\). Its normalized Yamabe functional is
\[
\mathcal{Y}(h)=\frac{\displaystyle\int_M R_h\,\dd\mu_h}{\displaystyle\operatorname{Vol}_h(M)^{(n-2)/n}}.
\]
Set
\[
Y_G(M,[g])=\inf_{h\in[g]_G}\mathcal{Y}(h).
\]
The question, under the stated conditional interpretation, is whether for every such \((M,g)\) and every subgroup \(G\leq\operatorname{Isom}(M,g)\), there exists \(h_*\in[g]_G\) satisfying
\[
\mathcal{Y}(h_*)=Y_G(M,[g]),
\qquad\text{equivalently}\qquad
\mathcal{Y}(h_*)\leq\mathcal{Y}(h)\quad\text{for every }h\in[g]_G.
\]
The supplied boundary-related reference does not specify a boundary version of this question. If manifolds with boundary, noncompact manifolds, or other dimensions are intended, the admissible class and functional require further specification.
\subsection{Short English statement}
Does every isometric group action admit a conformal metric that is invariant under the group and minimizes the normalized total scalar curvature among all invariant conformal metrics? The displayed formulation uses an explicitly conditional closed-manifold interpretation in dimension \(n\geq3\), because the source record omits the necessary geometric scope. The supplied partially solved label is dated curation, not an independently verified conclusion.
\subsection{Sources}
\begin{itemize}
\item Local UnsolvedMath registry, permanent ID \(30002635\), source problem code \texttt{catalogue entry}, title: \emph{Equivariant Yamabe Minimizers for Isometric Group Actions}; collection: \emph{Oberwolfach Reports Open Problems}. The supplied complete record is the local evidence used here. No individual proposer is identified.
\item Supplied report citation: \emph{Analysis, Geometry and Topology of Positive Scalar Curvature Metrics} (2015), Oberwolfach Reports, pp.\ 1991--2046, DOI: \url{https://doi.org/10.4171/owr/2014/36}. This is also the supplied literature link labelled ``original OWR report.'' The record separately gives proposed year 2014; the year discrepancy is retained without resolution.
\item Supplied literature reference labelled ``Madani equivariant Yamabe results'': \url{https://doi.org/10.1016/j.crma.2012.10.004}.
\item Supplied literature reference labelled ``equivariant Yamabe with boundary'': \url{https://doi.org/10.1007/s00526-021-02154-8}.
\item Literature assessment checked on \texttt{2026-08-21}: the supplied status is ``partially solved.'' The supplied assessment reports existence in broad classes described as nonpositive Yamabe cases, free actions, and proved low-dimensional or nondegenerate-action cases, while describing the fully general positive case as not established unconditionally. This is supplied dated curation derived from the OWR record, not independently verified literature assessment; the precise hypotheses and time-sensitive claims require independent review.
\end{itemize}
\end{document}
